% GENERATED FILE. Edit canonical lessons, then run npm run generate:books.
% canonical-source-hash: fnv1a64:e30f5a103ae809b5
% canonical-lessons: SA-C314-sauryam, SA-C314-dhirah, SA-C314-kirtih, SA-C314-gauravam, SA-C314-abhinandati

\chapter{A Brave Deed}
\label{ch:sa-a-brave-deed}
\hlchaptermodality{314}

\begin{chapteropening}
\textbf{By the end of this chapter:} \emph{I can tell of a brave deed and how it was honoured.}

Complete the last lesson of chapter 314: I can tell of a brave deed and how it was honoured.
\end{chapteropening}

\section[\texorpdfstring{\sk{शौर्यम्} (śauryam)}{शौर्यम् (śauryam)}]{\sk{शौर्यम्} (śauryam) — valour; heroism}

\label{lesson:SA-C314-sauryam}



\begin{quote}
Before the new one: say the Sanskrit for surrounds, then the Sanskrit for quells.
\end{quote}

\subsection*{You'll want to know: \sk{शौर्यम्}}
\textbf{\sk{शौर्यम्}} — \emph{śauryam} — \textquotedblleft{}valour; heroism\textquotedblright{}.

\begin{quote}
\sk{शौर्येण} \sk{सः} \sk{आहताय} \sk{औषधम्} \sk{आनयति।}
\end{quote}

\begin{quote}
\emph{śauryeṇa saḥ āhatāya auṣadham ānayati.}
\end{quote}

\begin{quote}
With great valour he brings medicine to the injured man.
\end{quote}

A neuter noun.

\begin{cousinweb}
\textbf{\sk{शौर्यम्}} (\emph{śauryam}) is usually traced to \emph{śūra}, \textquotedblleft{}a hero\textquotedblright{}, with the lengthened first vowel; MW gives \emph{śaurya} as heroism, valour, prowess (MW).
\end{cousinweb}

\begin{grammarlens}[title={What you've built}]
One more word for telling what happened.
\end{grammarlens}

\subsection*{Your turn}
\begin{itemize}
\raggedright
  \item \emph{Say it:} \emph{śauryam}
  \item \emph{Say it:} \emph{śauryam}, once more
  \item \emph{Say it:} \emph{śauryeṇa saḥ āhatāya auṣadham ānayati.} — With great valour he brings medicine to the injured man.
  \item \emph{Recall:} say the Sanskrit for surrounds, then the Sanskrit for quells, then say \emph{śauryam} again
\end{itemize}

\begin{tcolorbox}[breakable,colback=teal!4,colframe=teal!35!black,title={Before you move on}]
What does \sk{शौर्यम्} mean? (\textquotedblleft{}Valour; heroism\textquotedblright{}.) Say it once more.
\end{tcolorbox}
\section[\texorpdfstring{\sk{धीरः} (dhīraḥ)}{धीरः (dhīraḥ)}]{\sk{धीरः} (dhīraḥ) — steadfast; brave}

\label{lesson:SA-C314-dhirah}



\begin{quote}
Before the new one: say the Sanskrit for quells, then the Sanskrit for valour.
\end{quote}

\subsection*{You'll want to know: \sk{धीरः}}
\textbf{\sk{धीरः}} — \emph{dhīraḥ} — \textquotedblleft{}steadfast; brave\textquotedblright{}.

\begin{quote}
\sk{धीरः} \sk{सैनिकः} \sk{न} \sk{बिभेति।}
\end{quote}

\begin{quote}
\emph{dhīraḥ sainikaḥ na bibheti.}
\end{quote}

\begin{quote}
The steadfast soldier is not afraid.
\end{quote}

\begin{cousinweb}
\textbf{\sk{धीरः}} (\emph{dhīraḥ}) MW gives \emph{dhīra} as steady, firm, resolute, brave, and suggests, with a question mark, the root √dhṛ, \textquotedblleft{}to hold\textquotedblright{}, or √dhā, \textquotedblleft{}to put\textquotedblright{}, so its origin is uncertain (MW).
\end{cousinweb}

\begin{grammarlens}[title={What you've built}]
One more word for telling what happened.
\end{grammarlens}

\subsection*{Your turn}
\begin{itemize}
\raggedright
  \item \emph{Say it:} \emph{dhīraḥ}
  \item \emph{Say it:} \emph{dhīraḥ}, once more
  \item \emph{Say it:} \emph{dhīraḥ sainikaḥ na bibheti.} — The steadfast soldier is not afraid.
  \item \emph{Recall:} say the Sanskrit for quells, then the Sanskrit for valour, then say \emph{dhīraḥ} again
\end{itemize}

\begin{tcolorbox}[breakable,colback=teal!4,colframe=teal!35!black,title={Before you move on}]
What does \sk{धीरः} mean? (\textquotedblleft{}Steadfast; brave\textquotedblright{}.) Say it once more.
\end{tcolorbox}
\section[\texorpdfstring{\sk{कीर्तिः} (kīrtiḥ)}{कीर्तिः (kīrtiḥ)}]{\sk{कीर्तिः} (kīrtiḥ) — fame; renown}

\label{lesson:SA-C314-kirtih}



\begin{quote}
Before the new one: say the Sanskrit for valour, then the Sanskrit for steadfast.
\end{quote}

\subsection*{You'll want to know: \sk{कीर्तिः}}
\textbf{\sk{कीर्तिः}} — \emph{kīrtiḥ} — \textquotedblleft{}fame; renown\textquotedblright{}.

\begin{quote}
\sk{तस्य} \sk{कीर्तिः} \sk{सर्वत्र} \sk{प्रसरति।}
\end{quote}

\begin{quote}
\emph{tasya kīrtiḥ sarvatra prasarati.}
\end{quote}

\begin{quote}
His fame spreads everywhere.
\end{quote}

A feminine noun.

\begin{cousinweb}
\textbf{\sk{कीर्तिः}} (\emph{kīrtiḥ}) comes from the root that MW numbers √2. kṛ, \textquotedblleft{}to mention, praise\textquotedblright{}; MW gives \emph{kīrti} as mention, report and good report, fame (MW).
\end{cousinweb}

\begin{grammarlens}[title={What you've built}]
One more word for telling what happened.
\end{grammarlens}

\subsection*{Your turn}
\begin{itemize}
\raggedright
  \item \emph{Say it:} \emph{kīrtiḥ}
  \item \emph{Say it:} \emph{kīrtiḥ}, once more
  \item \emph{Say it:} \emph{tasya kīrtiḥ sarvatra prasarati.} — His fame spreads everywhere.
  \item \emph{Recall:} say the Sanskrit for valour, then the Sanskrit for steadfast, then say \emph{kīrtiḥ} again
\end{itemize}

\begin{tcolorbox}[breakable,colback=teal!4,colframe=teal!35!black,title={Before you move on}]
What does \sk{कीर्तिः} mean? (\textquotedblleft{}Fame; renown\textquotedblright{}.) Say it once more.
\end{tcolorbox}
\section[\texorpdfstring{\sk{गौरवम्} (gauravam)}{गौरवम् (gauravam)}]{\sk{गौरवम्} (gauravam) — honour; respect}

\label{lesson:SA-C314-gauravam}



\begin{quote}
Before the new one: say the Sanskrit for steadfast, then the Sanskrit for fame.
\end{quote}

\subsection*{You'll want to know: \sk{गौरवम्}}
\textbf{\sk{गौरवम्}} — \emph{gauravam} — \textquotedblleft{}honour; respect\textquotedblright{}.

\begin{quote}
\sk{ग्रामः} \sk{तस्मै} \sk{गौरवं} \sk{ददाति।}
\end{quote}

\begin{quote}
\emph{grāmaḥ tasmai gauravaṃ dadāti.}
\end{quote}

\begin{quote}
The village honours him.
\end{quote}

A neuter noun.

\begin{cousinweb}
\textbf{\sk{गौरवम्}} (\emph{gauravam}) is built on \emph{guru}, \textquotedblleft{}heavy, venerable\textquotedblright{}, with the lengthened first vowel; MW gives \emph{gaurava} as weight, venerableness, and respect shown to a person (MW).
\end{cousinweb}

\begin{grammarlens}[title={What you've built}]
One more word for telling what happened.
\end{grammarlens}

\subsection*{Your turn}
\begin{itemize}
\raggedright
  \item \emph{Say it:} \emph{gauravam}
  \item \emph{Say it:} \emph{gauravam}, once more
  \item \emph{Say it:} \emph{grāmaḥ tasmai gauravaṃ dadāti.} — The village honours him.
  \item \emph{Recall:} say the Sanskrit for steadfast, then the Sanskrit for fame, then say \emph{gauravam} again
\end{itemize}

\begin{tcolorbox}[breakable,colback=teal!4,colframe=teal!35!black,title={Before you move on}]
What does \sk{गौरवम्} mean? (\textquotedblleft{}Honour; respect\textquotedblright{}.) Say it once more.
\end{tcolorbox}
\section[\texorpdfstring{\sk{अभिनन्दति} (abhinandati)}{अभिनन्दति (abhinandati)}]{\sk{अभिनन्दति} (abhinandati) — congratulates; hails}

\label{lesson:SA-C314-abhinandati}



\begin{quote}
Before the new one: say the Sanskrit for fame, then the Sanskrit for honour.
\end{quote}

\subsection*{You'll want to know: \sk{अभिनन्दति}}
\textbf{\sk{अभिनन्दति}} — \emph{abhinandati} — \textquotedblleft{}congratulates; hails\textquotedblright{}.

\begin{quote}
\sk{राजा} \sk{वीरम्} \sk{अभिनन्दति।}
\end{quote}

\begin{quote}
\emph{rājā vīram abhinandati.}
\end{quote}

\begin{quote}
The king congratulates the hero.
\end{quote}

Past (he …): \sk{सः} \sk{वीरम्} \sk{अभिनन्दितवान्}, \emph{saḥ vīram abhinanditavān}.

\begin{cousinweb}
\textbf{\sk{अभिनन्दति}} (\emph{abhinandati}) is built from \emph{abhi} and the root √nand, \textquotedblleft{}to rejoice\textquotedblright{}, as in \emph{nandati}; MW gives \emph{abhi-nand} as to rejoice at, welcome, greet, hail, praise (MW).
\end{cousinweb}

\begin{grammarlens}[title={What you've built}]
One more word for telling what happened.
\end{grammarlens}

\subsection*{Your turn}
\begin{itemize}
\raggedright
  \item \emph{Say it:} \emph{abhinandati}
  \item \emph{Say it:} \emph{abhinandati}, once more
  \item \emph{Say it:} \emph{rājā vīram abhinandati.} — The king congratulates the hero.
  \item \emph{Recall:} say the Sanskrit for fame, then the Sanskrit for honour, then say \emph{abhinandati} again
\end{itemize}

\begin{tcolorbox}[breakable,colback=teal!4,colframe=teal!35!black,title={Before you move on}]
What does \sk{अभिनन्दति} mean? (\textquotedblleft{}Congratulates; hails\textquotedblright{}.) Say it once more.
\end{tcolorbox}
